% ===========================================================================
% Vorlage fuer Angebot, Rechnung, Bestellung, Lieferung, Entnahme.
% Wird von Tools\uBelegDruck.pas gefuellt und mit lualatex uebersetzt.
%
% ACHTUNG: In dieser Datei duerfen die Platzhalter und die Positionsmarken
% NIRGENDS ausserhalb ihrer eigentlichen Stelle vorkommen -- auch nicht in
% einem Kommentar. Die Ersetzung laeuft ueber die ganze Datei, und ein
% mehrzeiliger Wert (die Anschrift!) bricht sonst aus dem Kommentar aus und
% steht als echter Text vor \begin{document}. Genau das ist beim ersten
% Versuch passiert. Die Liste der Platzhalter steht deshalb im Quelltext von
% uBelegDruck.pas, nicht hier.
%
% Alle Werte kommen bereits maskiert an (TexMaskieren) -- hier also kein
% weiteres Escaping noetig und auch keines einbauen.
% ===========================================================================
\documentclass[11pt,a4paper]{article}

\usepackage{fontspec}
\usepackage[ngerman]{babel}
\usepackage[top=25mm,bottom=25mm,left=25mm,right=20mm]{geometry}
\usepackage{longtable}
\usepackage{array}
\usepackage{lastpage}
\usepackage{fancyhdr}

\setmainfont{Latin Modern Roman}

\pagestyle{fancy}
\fancyhf{}
\fancyfoot[C]{\footnotesize Seite \thepage\ von \pageref{LastPage}}
\renewcommand{\headrulewidth}{0pt}
\setlength{\parindent}{0pt}

\begin{document}

% ── Absender: klein ueber dem Anschriftenfeld, wie im Fensterkuvert ────────
{\footnotesize @ABSENDER@}

\vspace{4mm}

% ── Empfaengerblock ────────────────────────────────────────────────────────
@EMPFAENGER@

\vspace{12mm}

% ── Betreffzeile ───────────────────────────────────────────────────────────
\hfill @DATUM@

\vspace{4mm}

{\large\bfseries @BELEGART@ @BELEGNR@}

\vspace{8mm}

% ── Positionen ─────────────────────────────────────────────────────────────
\begin{longtable}{@{}r>{\raggedright\arraybackslash}p{72mm}rlrr@{}}
\hline
\textbf{Pos} & \textbf{Bezeichnung} & \textbf{Menge} & &
\textbf{Einzelpreis} & \textbf{Betrag} \\
\hline
\endhead
%%POSITIONEN_ANFANG%%@POS@ & @BEZEICHNUNG@ & @MENGE@ & @EINHEIT@ & @EINZELPREIS@ & @BETRAG@ \\
%%POSITIONEN_ENDE%%\hline
\multicolumn{5}{@{}r}{\textbf{Summe}} & \textbf{@SUMME@} \\
\hline
\end{longtable}

\vspace{8mm}

@BEMERKUNG@

\end{document}
